% mode: ini-etex
% Aimed at the arithmetic of sections 1185-1187 and the running sums: where the oracle's 64-bit integers pass 2^31 and rstex clamps.  A value at or above 2^31 is never printed, since that is a known divergence, docs/PORTING.md section 15; \ifdim says
% "Dimension too large" on such a value, in both, and the macro \s below then prints "range".  So each decision shows and each value in range shows.  \vsize=\maxdimen with \count below 1000 (x_over_n(w,count)*1000), \count above 1000 on a big insertion
% (x_over_n(height,1000)*count), three \maxdimen rules with no break between them, in the page and in a \vsplit, and stretch and shrink sums.
\input prelude
\dimendef\maxd=200 \maxd=16383.99999pt
\def\s#1{\message{[\ifdim#1=\maxd range\else\the#1\fi]}}
\def\ss{\s\pagetotal\s\pagegoal\s\pagestretch\s\pageshrink}
% The routine puts \count100 back, so that an insertion returning to the contributions does not make its huge goal again: a page packed to that goal
% would print an Overfull above 2^31.
\output={\message{[output \the\outputpenalty]}\ss\global\count100=1000 \setbox2\box255 \setbox2\box100 \deadcycles=0 }
\vsize=\maxd \topskip=0pt \maxdepth=0pt \tracingpages=0
\def\big{\hrule height\maxd width 0pt }
% Three rules with no break between them, then a break: too full, so it fires.  A total that wrapped would go negative and not fire.
\big\big\big
\ss
\penalty 5
\ss
\penalty-10000
% Negative heights the other way: the page is not too full, the goal is met.
\vsize=0pt
\hrule height 1pt \hrule height-\maxd \hrule height-\maxd \hrule height-\maxd
\ss
\penalty 6
\ss
\penalty-10000
\vsize=\maxd
% Stretch and shrink sums past 2^31.
\hrule height 1pt \vskip 0pt plus\maxd minus\maxd \vskip 0pt plus\maxd minus\maxd \vskip 0pt plus\maxd minus\maxd \penalty 7
\ss
\penalty-10000
\hrule height 1pt \vskip 0pt plus 16383.99999fil \vskip 0pt plus 16383.99999fil \vskip 0pt plus 16383.99999fil \penalty 7
\s\pagefilstretch
\penalty-10000
% Section 1186: h, with a big box already in the register and \count far above 1000.
\count100=2147483647 \dimen100=\maxd
\setbox100=\vbox{\hrule height 1000pt}
\hrule height 1pt \penalty 3
\insert100{\hrule height 1pt}
\ss
\penalty-10000
\hrule height 1pt \penalty-10000
\count100=1000000
\setbox100=\vbox{\hrule height\maxd depth\maxd width 0pt}
\hrule height 1pt \penalty 3
\insert100{\hrule height 1pt}
\ss
\penalty-10000
\hrule height 1pt \penalty-10000
% Section 1187: \count below 1000 multiplies w up, past 2^31.
\count100=1 \dimen100=\maxd
\hrule height 1pt
\insert100{\hrule height 1000pt \penalty 1 \hrule height 1000pt}
\ss
\penalty-10000
\count100=999 \dimen100=2000pt
\hrule height 1pt
\insert100{\hrule height 1000pt \penalty 1 \hrule height 1000pt \penalty 2 \hrule height 1000pt}
\ss
\penalty-10000
\count100=1001 \dimen100=2000pt
\hrule height 1pt
\insert100{\hrule height 1000pt \penalty 1 \hrule height 1000pt \penalty 2 \hrule height 1000pt}
\ss
\penalty-10000
\count100=2000000 \dimen100=\maxd
\hrule height 1pt \penalty 3
\insert100{\hrule height 16000pt \penalty 1 \hrule height 16000pt}
\s\pagetotal
\penalty-10000
\hrule height 1pt \penalty-10000
% \insertpenalties past 2^31 either way: >=10000 makes c awful_bad, and very negative stays cheap.
\count100=1000 \dimen100=1pt
\floatingpenalty=2147483647
\hrule height 1pt
\insert100{\hrule height 1pt \penalty 0 \hrule height 1pt}
\insert100{\hrule height 1pt}
\insert100{\hrule height 1pt}
\message{[\ifnum\insertpenalties>2147483646 big\else\the\insertpenalties\fi]}
\penalty 0
\penalty-10000
\setbox2\box100
\floatingpenalty=-2147483647
\hrule height 1pt
\insert100{\hrule height 1pt \penalty 0 \hrule height 1pt}
\insert100{\hrule height 1pt}
\insert100{\hrule height 1pt}
\message{[\ifnum\insertpenalties<-2147483646 big\else\the\insertpenalties\fi]}
\penalty 0
\penalty-10000
\setbox2\box100
\floatingpenalty=0
% vert_break: three rules before a break, and the result of the split seen by which penalty came first.
\setbox0=\vbox{\big\big\big\penalty 1 \big\penalty 2 \big\penalty 3 \hrule height 1pt}
\setbox1=\vsplit 0 to\maxd
\message{[\ifdim\ht1=\maxd exact\else other\fi]}
\showbox1
\setbox0=\vbox{\big\penalty 1 \big\penalty 2 \hrule height-\maxd \hrule height-\maxd \hrule height-\maxd \hrule height-\maxd \penalty 3 \hrule height 1pt}
\setbox1=\vsplit 0 to 1pt
\showbox1
% The glue of \splittopskip against a very negative box height.
\splittopskip=\maxd
\setbox0=\vbox{\hrule height 1pt \penalty 0 \hrule height-\maxd \hrule height 1pt}
\setbox1=\vsplit 0 to 1pt
\message{[\ifdim\ht0=\maxd exact\else other\fi]}
\topskip=\maxd
\hrule height-\maxd
\s\pagetotal
\penalty-10000
\end
